% "World of Adventure" campaign guide - Magic
% Written by Christopher Thomas.

\chapter{Magic}
\label{sect-magic}

\fixme{NYI}

note files:
- concentration
- constructs
- crafting dcs (alchemical transformations)
- magic background
- spells elements
- spells protect communal

To write up:
\begin{itemize}
\item Alchemical transmutations (base, precious, and perfect materials).
\item Notes about re-skinned spells, ``communal'' and ``mass'' spells, etc.
\item Notes about the Immaterial Realm and about divine magic.
\item Notes about magical background.
\end{itemize}

%
\section{The Nature of Magic}

\fixme{NYI}
%

%
\section{Spells Everyone Uses}

There are many spells (mostly cantrips and orisons) that see extremely
widespread use. These are listed below; full descriptions of modified or
nonstandard spells are in Section \ref{sect-spells}.

\begin{tabular}{ll}\hline
\textbf{Cantrips and Orisons} & \textbf{Notes} \\ \hline
Create Water & almost all drinking water is made with this \\
Detect Background & measure the level of background magic \\
Detect Magic & see active effects, see if an object is enchanted \\
Light & light bulb equivalent; most houses are lit with this \\
Mage Hand & telekinesis for lightweight unattended objects \\
Mending & make minor repairs to objects \\
Prestidigitation & clean stains, reflavour food, heat/chill drinks, etc \\
Purify Food and Drink & un-spoils food/water if iceboxes aren't available \\
Ray of Frost & touch-range variant is used for freezing food \\
Spark & disposable lighter equivalent; much better than flint/steel \\
Stabilize & stops someone at negative wounds from getting worse \\
\hline
\end{tabular}

\begin{tabular}{ll}\hline
\textbf{Other Spells} & \textbf{Notes} \\ \hline
Conjure Team & virtually all draft animals are conjured \\
Diagnose Disease & medical diagnostics spell \\
Diagnose Injury & medical diagnostics spell \\
Diagnose Poison & medical diagnostics spell \\
Floating Disk & lightweight cargo transport \\
Freight Pallet & heavy cargo transport \\
Mount & virtually all riding mounts are conjured \\
P from Background & explore or live in weakly magical areas safely \\
P from Mental Control & public officials usually have this active \\
Public Address & microphone and loundspeaker equivalent \\
\hline
\end{tabular}
%

%
\section{Magic Items Everyone Uses}

There are many types of magic item that are very commonly made and that are
within reach of individuals, and there are several more that are more
expensive but that are available for use (either for a fee or as public
infrastructure).

One-shot items:

\begin{tabular}{ll} \hline
\textbf{Item} & \textbf{Notes} \\ \hline
Knock & police breaching tool \\
Neutralize Poison & expensive field treatment \\
\hline
\end{tabular}

Personal enchantments (part of an amulet, chain of office, cloak, hat, or
ring):

\begin{tabular}{ll} \hline
\textbf{Item} & \textbf{Notes} \\ \hline
Comprehend Languages & ring, amulet, or hat \\
Disguise Self & hat or clothing; police investigator tool \\
Endure Elements & cloak or amulet \\
P from Background & amulet \\
P from Mental Control & amulet, chain of office \\
Speak with Animals & ring, amulet, or circlet; druid / vet tool \\
\hline
Diplomacy bonus & used by politicians; also Bluff and Sense Motive \\
Prof/Craft bonus & used by most late-career professionals \\
\hline
\end{tabular}

Public infrastructure enchantments (usually wondrous items):

\begin{tabular}{ll} \hline
\textbf{Item} & \textbf{Notes} \\ \hline
Contraception & cantrip; devices installed in public locations \\
Create Food & enabled transition from agrarian to manufacturing society \\
Scrying & crystal balls are used by law enforcement \\
Sending & magic mirrors are owned by the wealthy \\
\hline
\end{tabular}
%

%
\section{Magic Item Prices}

Magic item prices are calculated per the CRB's rules, with the following
changes:
\begin{itemize}
%
\item Prices for \textit{spell effects} are on an exponential scale rather
than a quadratic scale, as described below.
%
\item For \textbf{standard pricing}, prices are in silver pieces rather than
in gold pieces. Magic items are still quite expensive but are in reach of
late-career professionals. It is cheaper to hire a spellcaster on salary for
utility spells than to get an enchanted item.
%
\item For \textbf{cheaper spell magic}, prices for \textit{spell effects}
are reduced by another factor of 10. This puts enchanted items in the
``expensive appliance'' price range; spellcasters are the ones crafting the
items, rather than being hired to cast spells.
%
\end{itemize}

Notes about designing magic items:
\begin{itemize}
%
\item If there are several ways to get an effect, only the more expensive
one works. For example, a continuous Cat's Grace item does not work - build
an ability boost item instead. Similarly, a continuous item of Mage Armour
does not work - make armour with an enhancement bonus instead.
%
\item Wands are always charged (no unlimited-charge wands). Wands that run
out of charges become nonmagical, and have their enchantable materials
degraded to unusable base forms, but otherwise remain intact. Likewise with
other charged or one-shot items.
%
\item Unlimited-use items tend to be wondrous items shaped like appropriate
tools. Unlimited-use items \textbf{only operate at touch range}.
%
\item Items that are slotted should have effects that are appropriate to the
slot. A Hat of Diplomacy or Vambraces of Natural Armour are reasonable.
Vambraces of Diplomacy aren't. This is adjudicated on a case-by-case basis
by the DM.
%
\item Items that have material component costs may either have the component
cost built in (cost multiplied by number of charges), or may be built to
consume externally-supplied components (no added cost, but expending charges
or uses without components wastes those uses).
%
\item Enchanting a spell effect (or other effect that has a spell
requirement) requires either having that spell or having the
\textit{full-time help} of someone who has that spell. At the DM's
discretion, the requirement can be waived if the crafting DC is increased
(per the SRD), but the crafter must still \textit{qualify for being able to
get that spell}.
%
\item Enchanting an armour or weapon enhancement bonus requires the crafter
to have a BAB equal to or greater than the bonus, or to have the
\textit{full-time help} of someone who has that BAB.
%
\item Enchanting a skill bonus requires the crafter to have a number of
ranks in that skill equal to or greater than the bonus, or to have the
\textit{full-time help} of someone who has that many ranks.
%
\item Other types of bonus have similar prerequisites, adjudicated by the
DM on a case-by-case basis.
%
\end{itemize}

The CRB's prices for spell effects scale as \textbf{spell level $\times$
caster level}. The house rules \textbf{replace this} with the following
scale:

\begin{tabular}{|l|l|lll|lll|lll|} \hline
\textbf{Spell Level} & 0 & 1 & 2 & 3 & 4 & 5 & 6 & 7 & 8 & 9 \\ \hline
\textbf{Price Factor} & $\frac{1}{2}$ & 1 & 3 & 10 & 30 & 100 & 300 &
1000 & 3000 & 10000 \\ \hline
\textbf{Caster Level} & 1 & 1 & 3 & 5 & 7 & 9 & 11 & 13 & 15 & 17 \\ \hline
\end{tabular}

Spells can be enchanted with higher caster level. When doing so, scale the
price by \textbf{new caster level $\div$ minimum caster level}.

The net effect of the scale is as follows:
\begin{itemize}
\item Multipliers for L1 and L4 spells are about the same (1 and 30 vs 1 and
28).
\item L2 and L3 items are cheaper. In particular there is no longer a 6x
price jump from SL1 to SL2.
\item Effects above L4 get very expensive very fast. A Ring of Wishing is
astronomically expensive.
\end{itemize}

The base price of spell effects in items is as follows (adapted from the SRD):

\begin{tabular}{llll} \hline
\textbf{Type} & \textbf{Example} & \textbf{Standard} & \textbf{Even Cheaper}
\\ \hline
Spell Completion & scroll & PF $\times$ 25 sp & PF $\times$ 2.5 sp \\
One-Shot Use & potion & PF $\times$ 50 sp & PF $\times$ 5 sp \\
Spell Trigger 50ch & wand & PF $\times$ 750 sp & PF $\times$ 75 sp \\
Command Word & wondrous & PF $\times$ 1800 sp & PF $\times$ 180 sp \\
Use or Continuous & wondrous & PF $\times$ 2000 sp & PF $\times$ 200 sp \\
\hline
\end{tabular}

The base price of other effects in items is as follows (adapted from the SRD):

\begin{tabular}{lll} \hline
\textbf{Effect} & \textbf{Base Price} & \textbf{Notes} \\ \hline
Ability Bonus (enhance) & bonus$^2 \times$ 1000 sp & \\
Armour Bonus (enhance) & bonus$^2 \times$ 500 sp &
applies to armour rating, not to AC or DR \\
AC Bonus (deflect) & bonus$^2 \times$ 2000 sp & applies directly to AC \\
AC Bonus (other) & bonus$^2 \times$ 2500 sp & applies directly to AC \\
Natural Armour Bonus (enhance) & bonus$^2 \times$ 1000 sp &
applies to armour rating, not to AC or DR \\
Save Bonus (resist) & bonus$^2 \times$ 1000 sp & \\
Save Bonus (other) & bonus$^2 \times$ 2000 sp & \\
Skill Bonus (competence) & bonus$^2 \times$ 100 sp & \\
Weapon Bonus (enhance) & bonus$^2 \times$ 2000 sp &
applies to to-hit \textit{and} damage \\
\hline
\end{tabular}

I don't use Spell Resistance, so that effect isn't in the list.

Per Section \ref{sect-equip}, enchantments take up ``slots'' in items.
Certain other effects (such as making an item into a wizard's bonded object
or a cleric's holy symbol) also take up slots, but have no additional cost
beyond the cost of making the item enchantable.
%

%
\section{Inventing New Spells}

Inventing new spells is treated as an extended crafting project, using the
Spellcraft skill. The target DC and the project cost depend on the target
spell's level and on how novel the spell is.


Reinventing a spell that is widely known, or a simple variant of such
a spell:

\mbox{\bfseries DC 15 + (2 $\times$ spell level)}, cost
\mbox{\bfseries PF $\times$ 1000 sp}

Inventing a spell that is novel but that fits within established magical
paradigms:

\mbox{\bfseries DC 20 + (2 $\times$ spell level)}, cost
\mbox{\bfseries PF $\times$ 2000 sp}

Inventing a spell that breaks new ground and is the first of its kind:

\mbox{\bfseries DC 25 + (2 $\times$ spell level)}, cost
\mbox{\bfseries PF $\times$ 4000 sp}

The ``materials and supplies cost'' of the project is the entire value, not
half the value, for this sort of research.
%

%
\section{Special Materials}

\fixme{NYI}
%

%
\section{Communal and Mass Spells}

\fixme{NYI}
%

%
\section{Re-Skinned Combat Spells}

\fixme{NYI}
%

%
\section{Constructs}

\fixme{NYI}

%
% This is the end of the file.
